\solchapter{cosmo}

\sol{cosmo}{1} $1/H_0 = 4.4\times10^{17}\un{s} = 14$ billion years. A galaxy at $100\un{Mpc}$ recedes at $7000\un{km/s}$. The limit $c/H_0 = 4300\un{Mpc}$ is the Hubble radius. Relativity forbids anything from overtaking a light ray \emph{locally}; it says nothing about how fast the distance between distant objects can grow when space itself expands. Nothing ever moves past light in its own neighbourhood, so no signal ever outruns light.

\sol{cosmo}{2} $\rho_c = 3H_0^2/8\pi G = 9.2\times10^{-27}\un{kg/m^3}$, or $5.5$ protons per cubic metre. Air ($1.2\un{kg/m^3}$) is $10^{26}$ times denser, and even the interstellar gas of our galaxy is $2\times10^5$ times denser. The universe is mostly emptier than the best laboratory vacuum.

\sol{cosmo}{3} $T = 2.725\un{K}\times1101 = 3000\un{K}$. By Wien's law it peaked at $\lambda = (2.9\un{mm\,K})/T = 970\un{nm}$, in the near-infrared; today the peak is at $1.06\un{mm}$, in the microwave. A $3000\un{K}$ blackbody glows a dull orange, like a candle flame or an old incandescent bulb, and the whole sky would have been that colour.

\sol{cosmo}{4} $a\propto t^{2/3}$ gives $H = \dot a/a = 2/3t$, so $t_0 = 2/3H_0 = 9.3$ billion years. That is younger than the oldest stars: a matter-only universe is ruled out. With dark energy, the expansion was slower in the past than now, so for the same $H_0$ the universe is older, $13.8$ billion years.

\sol{cosmo}{5} By Newton's shell theorem, only the mass $M$ inside affects a galaxy on the sphere, so its energy per unit mass is constant:
\[
  \tfrac12\dot R^2 - \frac{GM}{R} = E.
\]
With $R = ar$ and $M = \tfrac43\pi a^3r^3\rho$, this is $\tfrac12\dot a^2r^2 - \tfrac43\pi G\rho a^2r^2 = E$. Divide by $\tfrac12a^2r^2$:
\[
  \left(\frac{\dot a}{a}\right)^2 = \frac{8\pi G}{3}\rho + \frac{2E}{r^2a^2}.
\]
Since the expansion is uniform, every shell must give the same equation, so $2E/r^2$ is the same constant for all of them; call it $-kc^2$. This is the Friedmann equation. Positive energy ($k<0$) means expansion forever; negative energy ($k>0$) means the universe is bound, and eventually recollapses (without dark energy). General relativity reinterprets $k$ as the curvature of space.

\sol{cosmo}{6} $\rho_r\propto a^{-4}$ and $\rho_m\propto a^{-3}$, so the ratio $\rho_r/\rho_m$ grows as $1/a = 1+z$ going back in time. They were equal at $1 + z = 0.31/(9\times10^{-5}) = 3400$. Integrating the Friedmann equation gives an age of about $50\,000$ years.

\sol{cosmo}{7} $n_b/n_\gamma = 0.25\un{m^{-3}}/(4.11\times10^8\un{m^{-3}}) = 6\times10^{-10}$. This tiny number is what remains of the slight excess of matter over antimatter in the early universe: for every billion antiprotons there were a billion and one protons, and after annihilation the one survived. It also controls much of later history. It sets the abundances of deuterium and helium made in the first minutes, and the enormous number of photons per baryon delays recombination until $3000\un{K}$, far below the $1.6\times10^5\un{K}$ that corresponds to hydrogen's $13.6\un{eV}$.

\sol{cosmo}{8} $3ct = 1.1\times10^{6}$ light-years at recombination, stretched to $1.3\times10^9$ light-years today. At $46$ billion light-years, it subtends $0.027\un{rad}\approx1.6^\circ$, about three times the width of the full Moon. In a standard Big Bang, patches of sky more than a couple of degrees apart were never in causal contact, yet they have the same temperature to one part in $10^5$. Nothing could have evened them out. Inflation solves the problem by having the whole observable universe emerge from a single tiny patch that was in contact before inflation stretched it.

\sol{cosmo}{9} At freeze-out, $n/p = e^{-1.29/0.8} = 0.20$, about $1/5$. After $180\un{s}$ of decay, it is $0.20\,e^{-180/880} = 0.16$. Every helium-4 nucleus contains $2$ neutrons and $2$ protons, so if all the neutrons end up in helium, the mass in helium is $2n$ out of $n+p$ nucleons: $Y = 2n/(n+p) = 2(n/p)/(1+n/p) = 0.28$. This is close to the observed $0.245$; a careful treatment, with gradual freeze-out, gives $n/p\approx1/7$ and $Y\approx0.25$. A two-line calculation from the weak interaction, the neutron lifetime, and the expansion rate predicts the composition of the universe.

\sol{cosmo}{10} From $a = 10^{-28}$ to matter--radiation equality at $a\approx3\times10^{-4}$, $\abs{\Omega-1}$ grows by $(3\times10^{24})^2\approx10^{49}$; from then to now, by another $3000$. So $\abs{\Omega-1}$ must have been less than $0.01/(3\times10^{52})\approx3\times10^{-55}$ at the time: fine-tuning to fifty-five decimal places. (Dark energy has recently begun to reduce $\abs{\Omega-1}$, but only slightly.) If inflation starts from $\abs{\Omega-1}\sim1$, it must shrink it by $e^{-2N}\lesssim3\times10^{-55}$, so $N\gtrsim\tfrac12\ln(3\times10^{54})\approx63$. Inflation models are designed to give at least $60$ e-folds, and it costs them little to give many more.

\sol{cosmo}{11} Stars of radius $R$ and number density $n$ in a shell of thickness $\dd r$ at distance $r$ cover a fraction $n\,4\pi r^2\dd r\cdot\pi R^2/4\pi r^2 = n\pi R^2\dd r$ of the sky, independent of $r$: farther shells contain more stars, each looking smaller. Adding infinitely many shells covers the sky, and a typical line of sight ends on a star after a distance $\ell = 1/n\pi R^2$. A star's surface brightness does not depend on its distance, so the sky would be uniformly as bright as the Sun's surface. With about $10^9$ stars per cubic megaparsec on average, $\ell\sim10^{24}$ light-years. The universe is only $14$ billion years old, so we see stars only out to about $10^{10}$ light-years, and they cover only about $10^{-14}$ of the sky. (And stars live only billions of years, so they could never fill the sky with light even given time.) The finite age matters far more; expansion dims distant starlight by redshift, but only by a modest factor. The dark night sky tells you the universe had a beginning.
